% ============================================================
% CHAPITRE 5 — SPRINT 2 : INTELLIGENCE ARTIFICIELLE ET ANALYTICS
% ============================================================

\chapter{Sprint 2 --- Intelligence Artificielle et Analytics}

\section*{Plan}

\begin{tabular}{r l r}
\textbf{1} & \textbf{Sprint Planning} & \textbf{\pageref{sec:s2_planning}} \\
\textbf{2} & \textbf{Analyse des User Stories} & \textbf{\pageref{sec:s2_stories}} \\
\textbf{3} & \textbf{Conception} & \textbf{\pageref{sec:s2_conception}} \\
\textbf{4} & \textbf{Réalisation et tests} & \textbf{\pageref{sec:s2_realisation}} \\
\textbf{5} & \textbf{Sprint Rétrospective} & \textbf{\pageref{sec:s2_retro}} \\
\end{tabular}

\section*{Introduction}

Le Sprint~2 constitue le c\oe ur intelligent de la plateforme LLM Dashboard. Il est consacré au développement des composants d'apprentissage automatique et d'analyse avancée~: la détection d'anomalies de consommation par trois algorithmes complémentaires, l'explication automatique des anomalies par intelligence artificielle, la prévision budgétaire à trois scénarios, l'optimisation des coûts et le pipeline d'agrégation des données via Celery. Ce sprint transforme la plateforme d'un simple outil de monitoring en un système \textbf{prédictif et prescriptif}.

% ────────────────────────────────────────────────────────────
\section{Sprint Planning}
\label{sec:s2_planning}
% ────────────────────────────────────────────────────────────

\begin{table}[H]
\centering
\caption{Sprint Backlog --- Sprint 2}
\label{tab:sprint2_backlog}
\begin{tabularx}{\textwidth}{|c|l|X|c|}
\hline
\textbf{ID} & \textbf{Module} & \textbf{User Story} & \textbf{Priorité} \\
\hline
US-7 & Detection & Détecter les anomalies de consommation via ML & Must \\
\hline
US-8 & Explainer & Expliquer automatiquement les anomalies par IA & Should \\
\hline
US-9 & Forecasting & Prévoir les coûts futurs selon trois scénarios & Should \\
\hline
US-10 & Analytics & Agréger et analyser les coûts par modèle et par jour & Must \\
\hline
US-11 & Optimizer & Recommander des optimisations de coûts & Should \\
\hline
\end{tabularx}
\end{table}


% ────────────────────────────────────────────────────────────
\section{Analyse des User Stories}
\label{sec:s2_stories}
% ────────────────────────────────────────────────────────────

\subsection{User Story~: Détection d'anomalies (US-7)}

\begin{table}[H]
\centering
\caption{User Story US-7~: Détection d'anomalies}
\label{tab:us7}
\begin{tabularx}{\textwidth}{|l|X|}
\hline
\textbf{Titre} & Détecter automatiquement les anomalies de consommation \\
\hline
\textbf{En tant que} & Administrateur tenant \\
\hline
\textbf{Je veux} & Être alerté lorsqu'une consommation anormale de tokens est détectée \\
\hline
\textbf{Afin de} & Réagir rapidement aux dérives de coûts et aux usages non conformes \\
\hline
\textbf{Critères d'acceptation} & \begin{minipage}[t]{\linewidth}\vspace{2pt}
\begin{itemize}[noitemsep,topsep=0pt]
\item Trois algorithmes ML complémentaires~: STL decomposition~\cite{cleveland1990}, Isolation Forest~\cite{liu2008}, CUSUM~\cite{page1954}
\item Baselines dynamiques par tenant, modèle et agent
\item Classification par sévérité~: info, warning, critical
\item Publication d'un événement \texttt{anomaly.detected} sur le bus d'événements
\item Stockage des anomalies avec métadonnées complètes dans la table \texttt{llm\_anomaly}
\end{itemize}\vspace{2pt}
\end{minipage} \\
\hline
\end{tabularx}
\end{table}

\subsection{User Story~: Explication IA des anomalies (US-8)}

\begin{table}[H]
\centering
\caption{User Story US-8~: Explication IA}
\label{tab:us8}
\begin{tabularx}{\textwidth}{|l|X|}
\hline
\textbf{Titre} & Comprendre les causes d'une anomalie grâce à l'IA \\
\hline
\textbf{En tant que} & Manager / Administrateur \\
\hline
\textbf{Je veux} & Recevoir une explication en langage naturel pour chaque anomalie \\
\hline
\textbf{Afin de} & Comprendre les causes racines et agir en conséquence \\
\hline
\textbf{Critères d'acceptation} & \begin{minipage}[t]{\linewidth}\vspace{2pt}
\begin{itemize}[noitemsep,topsep=0pt]
\item Génération automatique d'une explication structurée (contexte, cause, impact)
\item Support multi-fournisseur LLM (Groq, OpenAI, Anthropic, Gemini, Ollama)
\item Production de 2--4 recommandations actionnables classées par impact
\item Stockage des embeddings vectoriels (384 dimensions) via pgvector
\item Cache des explications dans Redis pour éviter les appels LLM redondants
\end{itemize}\vspace{2pt}
\end{minipage} \\
\hline
\end{tabularx}
\end{table}

\subsection{User Story~: Prévision budgétaire (US-9)}

\begin{table}[H]
\centering
\caption{User Story US-9~: Prévision budgétaire}
\label{tab:us9}
\begin{tabularx}{\textwidth}{|l|X|}
\hline
\textbf{Titre} & Prévoir les coûts futurs pour planifier le budget \\
\hline
\textbf{En tant que} & Décideur / Manager \\
\hline
\textbf{Je veux} & Visualiser les prévisions de coûts à 30, 60 et 90 jours \\
\hline
\textbf{Afin de} & Anticiper les dépenses et éviter les dépassements budgétaires \\
\hline
\textbf{Critères d'acceptation} & \begin{minipage}[t]{\linewidth}\vspace{2pt}
\begin{itemize}[noitemsep,topsep=0pt]
\item Trois scénarios de prévision~: optimiste, moyen, pessimiste~\cite{hyndman2018}
\item Calcul de la tendance quotidienne, du taux de croissance et de la saisonnalité
\item Identification des risques de dépassement avec alertes
\item Visualisation graphique des intervalles de confiance
\end{itemize}\vspace{2pt}
\end{minipage} \\
\hline
\end{tabularx}
\end{table}

\subsection{User Story~: Analytics et agrégation (US-10)}

\begin{table}[H]
\centering
\caption{User Story US-10~: Analytics}
\label{tab:us10}
\begin{tabularx}{\textwidth}{|l|X|}
\hline
\textbf{Titre} & Consulter les coûts agrégés par modèle et par jour \\
\hline
\textbf{En tant que} & Administrateur \\
\hline
\textbf{Je veux} & Visualiser les coûts agrégés quotidiennement et mensuellement \\
\hline
\textbf{Afin de} & Comprendre la répartition des dépenses LLM \\
\hline
\textbf{Critères d'acceptation} & \begin{minipage}[t]{\linewidth}\vspace{2pt}
\begin{itemize}[noitemsep,topsep=0pt]
\item Agrégation nocturne automatique via Celery Beat~\cite{celery2024}
\item Rollup mensuel pour alimenter la facturation
\item Pricing versionné avec plages de dates (table \texttt{llm\_model\_pricing})
\item Endpoint API \texttt{/v1/analytics/costs} avec filtres par période, modèle et provider
\end{itemize}\vspace{2pt}
\end{minipage} \\
\hline
\end{tabularx}
\end{table}

\subsection{User Story~: Optimisation (US-11)}

\begin{table}[H]
\centering
\caption{User Story US-11~: Optimisation des coûts}
\label{tab:us11}
\begin{tabularx}{\textwidth}{|l|X|}
\hline
\textbf{Titre} & Recevoir des recommandations d'optimisation \\
\hline
\textbf{En tant que} & Administrateur tenant \\
\hline
\textbf{Je veux} & Obtenir des recommandations pour réduire mes coûts LLM \\
\hline
\textbf{Afin de} & Optimiser mon budget sans compromettre la qualité de service \\
\hline
\textbf{Critères d'acceptation} & \begin{minipage}[t]{\linewidth}\vspace{2pt}
\begin{itemize}[noitemsep,topsep=0pt]
\item Recommandations de substitution de modèles avec simulation ROI
\item Recommandations d'optimisation structurelle de prompts
\item Cycle de vie~: proposée $\rightarrow$ validée $\rightarrow$ déployée / rejetée
\item Scan hebdomadaire automatique via Celery Beat
\end{itemize}\vspace{2pt}
\end{minipage} \\
\hline
\end{tabularx}
\end{table}


% ────────────────────────────────────────────────────────────
\section{Conception}
\label{sec:s2_conception}
% ────────────────────────────────────────────────────────────

\subsection{Diagramme de classes --- Module Detection}

\begin{figure}[H]
    \centering
    \fbox{\parbox{14cm}{\centering\vspace{4cm}\textit{Diagramme de classes --- Module Detection}\\[0.5cm]\textit{(Classes~: LLMAnomaly, LLMTokenBaseline, AnomalyDetectionService)}\\[0.5cm]\textit{(Attributs~: anomaly\_type, severity, z\_score, algorithm\_used, details JSONB)}\vspace{3cm}}}
    \caption{Diagramme de classes du module Detection}
    \label{fig:class_detection}
\end{figure}

\subsection{Diagramme de classes --- Modules Analytics, Forecasting et Optimizer}

Le diagramme de classes suivant modélise les entités centrales des modules d'agrégation de coûts, de prévision budgétaire et d'optimisation~:

\begin{figure}[H]
    \centering
    \fbox{\parbox{14cm}{\centering\vspace{4cm}\textit{Diagramme de classes --- Modules Analytics, Forecasting et Optimizer}\\[0.5cm]\textit{(Classes~: LLMCostDaily, LLMCostMonthly, LLMModelPricing, AnalyticsService,}\\
    \textit{ForecastResult, ForecastService, OptimizationRecommendation, OptimizerService)}\\[0.5cm]\textit{(Voir fichier PlantUML~: ch5\_class\_analytics.puml)}\vspace{3cm}}}
    \caption{Diagramme de classes des modules Analytics, Forecasting et Optimizer}
    \label{fig:class_analytics}
\end{figure}

\subsection{Diagramme de séquence --- Pipeline de détection d'anomalies}

\begin{figure}[H]
    \centering
    \fbox{\parbox{14cm}{\centering\vspace{4cm}\textit{Diagramme de séquence --- Pipeline de détection d'anomalies}\\[0.5cm]\textit{(EventBus $\rightarrow$ DetectionService $\rightarrow$ STL $\rightarrow$ IsolationForest $\rightarrow$ CUSUM)}\\[0.5cm]\textit{(1. Événement reçu \quad 2. Requête baseline \quad 3. Exécution des 3 algorithmes)}\\
    \textit{(4. Classification de sévérité \quad 5. Persistance \quad 6. Publication événement)}\vspace{3cm}}}
    \caption{Diagramme de séquence du pipeline de détection d'anomalies}
    \label{fig:seq_detection}
\end{figure}

\subsection{Diagramme de séquence --- Explication IA}

\begin{figure}[H]
    \centering
    \fbox{\parbox{14cm}{\centering\vspace{4cm}\textit{Diagramme de séquence --- Explication IA d'une anomalie}\\[0.5cm]\textit{(1. Événement anomaly.detected reçu \quad 2. Construction du prompt contextuel)}\\[0.5cm]\textit{(3. Appel LLM (Groq/OpenAI) \quad 4. Parsing de la réponse structurée)}\\
    \textit{(5. Génération des embeddings \quad 6. Stockage dans pgvector)}\vspace{3cm}}}
    \caption{Diagramme de séquence de l'explication IA}
    \label{fig:seq_explainer}
\end{figure}

\subsection{Diagramme de séquence --- Prévision budgétaire}

Le diagramme suivant détaille le flux de prévision budgétaire à trois scénarios (optimiste, moyen, pessimiste), incluant l'évaluation des risques de dépassement~:

\begin{figure}[H]
    \centering
    \fbox{\parbox{14cm}{\centering\vspace{4cm}\textit{Diagramme de séquence --- Prévision budgétaire à 3 scénarios}\\[0.5cm]\textit{(Admin $\rightarrow$ Frontend $\rightarrow$ API $\rightarrow$ ForecastService $\rightarrow$ DB $\rightarrow$ RiskAssessor)}\\[0.5cm]\textit{(1. Récupération des coûts sur 90j \quad 2. Calcul de tendance et croissance)}\\
    \textit{(3. Projection 3 scénarios \quad 4. Évaluation des risques \quad 5. Persistance)}\\[0.3cm]\textit{(Voir fichier PlantUML~: ch5\_seq\_forecast.puml)}\vspace{3cm}}}
    \caption{Diagramme de séquence de la prévision budgétaire}
    \label{fig:seq_forecast}
\end{figure}

\subsection{Diagramme de séquence --- Agrégation nocturne des coûts}

Le pipeline d'agrégation, orchestré par Celery Beat, s'exécute automatiquement chaque nuit et chaque premier du mois pour produire les rollups nécessaires à la facturation~:

\begin{figure}[H]
    \centering
    \fbox{\parbox{14cm}{\centering\vspace{4cm}\textit{Diagramme de séquence --- Pipeline d'agrégation Celery}\\[0.5cm]\textit{(Celery Beat $\rightarrow$ Worker $\rightarrow$ AnalyticsService $\rightarrow$ llm\_token\_log $\rightarrow$ llm\_cost\_daily $\rightarrow$ llm\_cost\_monthly)}\\[0.5cm]\textit{(1. Agrégation quotidienne \quad 2. Déclenchement détection anomalies)}\\
    \textit{(3. Rollup mensuel \quad 4. Publication événement pour facturation)}\\[0.3cm]\textit{(Voir fichier PlantUML~: ch5\_seq\_aggregation.puml)}\vspace{3cm}}}
    \caption{Diagramme de séquence du pipeline d'agrégation nocturne}
    \label{fig:seq_aggregation}
\end{figure}

\subsection{Diagramme d'activité --- Pipeline Celery}

\begin{figure}[H]
    \centering
    \fbox{\parbox{14cm}{\centering\vspace{4cm}\textit{Diagramme d'activité --- Pipeline d'agrégation Celery}\\[0.5cm]\textit{(Tâches périodiques orchestrées par Celery Beat~:)}\\[0.3cm]\textit{- Chaque nuit~: aggregate\_daily\_costs}\\
    \textit{- Chaque 1\textsuperscript{er} du mois~: rollup\_monthly\_costs}\\
    \textit{- Chaque semaine~: run\_optimization\_scan}\vspace{3cm}}}
    \caption{Diagramme d'activité du pipeline Celery}
    \label{fig:activity_celery}
\end{figure}


% ────────────────────────────────────────────────────────────
\section{Réalisation et tests}
\label{sec:s2_realisation}
% ────────────────────────────────────────────────────────────

\subsection{Module Detection --- Implémentation}

Le module de détection d'anomalies implémente les trois algorithmes présentés dans l'état de l'art (chapitre~2). Le service principal orchestre l'exécution séquentielle des détecteurs~:

\begin{lstlisting}[language=Python3, caption={Service de détection d'anomalies (modules/detection/services/detection\_service.py)}]
class AnomalyDetectionService:
    async def detect(self, tenant_id: str, 
                      lookback_days: int = 14) -> list[LLMAnomaly]:
        # 1. Recuperer les donnees historiques
        logs = await self.repo.get_token_logs(tenant_id, lookback_days)
        baseline = await self.repo.get_baseline(tenant_id)
        
        anomalies = []
        
        # 2. STL decomposition — detecte les deviations saisonnieres
        stl_results = self._run_stl(logs, baseline)
        anomalies.extend(stl_results)
        
        # 3. Isolation Forest — detecte les points aberrants
        if_results = self._run_isolation_forest(logs)
        anomalies.extend(if_results)
        
        # 4. CUSUM — detecte les derives progressives
        cusum_results = self._run_cusum(logs, baseline)
        anomalies.extend(cusum_results)
        
        # 5. Classifier par severite et dedupliquer
        classified = self._classify_severity(anomalies)
        
        # 6. Publier les evenements
        for anomaly in classified:
            await publish("anomaly.detected", anomaly.to_dict())
        
        return classified
\end{lstlisting}


\subsection{Module Explainer --- Implémentation}

Le module Explainer utilise un LLM pour générer des explications structurées des anomalies détectées. Un prompt contextuel est construit dynamiquement avec les données de l'anomalie~:

\begin{lstlisting}[language=Python3, caption={Service d'explication IA (modules/detection/services/explainer\_service.py)}]
class ExplainerService:
    async def explain(self, anomaly: LLMAnomaly) -> Explanation:
        # Construire le prompt contextuel
        prompt = self._build_prompt(anomaly)
        
        # Appeler le LLM (fallback automatique entre providers)
        for provider in [GroqProvider, OpenAIProvider, AnthropicProvider]:
            try:
                response = await provider.generate(prompt)
                break
            except ProviderError:
                continue
        
        # Parser la reponse structuree
        explanation = self._parse_explanation(response)
        
        # Generer l'embedding pour la recherche semantique
        embedding = self.embedder.encode(explanation.summary)
        
        # Stocker dans pgvector
        await self.repo.save(explanation, embedding)
        
        return explanation
\end{lstlisting}


\subsection{Module Forecasting --- Implémentation}

Le module de prévision budgétaire calcule trois scénarios (optimiste, moyen, pessimiste) fondés sur l'analyse de la tendance et de la saisonnalité des données historiques~\cite{hyndman2018}~:

\begin{lstlisting}[language=Python3, caption={Service de prévision budgétaire (modules/forecasting/services/forecast\_service.py)}]
class ForecastService:
    async def forecast(self, tenant_id: str, 
                        horizon_days: int = 30) -> ForecastResult:
        # Recuperer les couts quotidiens historiques
        daily_costs = await self.repo.get_daily_costs(tenant_id, days=90)
        
        # Calculer la tendance et le taux de croissance
        trend = self._compute_trend(daily_costs)
        growth_rate = self._compute_growth_rate(daily_costs)
        
        # Generer les 3 scenarios
        optimistic = self._project(trend, growth_rate * 0.5, horizon_days)
        baseline   = self._project(trend, growth_rate,       horizon_days)
        pessimistic = self._project(trend, growth_rate * 1.5, horizon_days)
        
        # Evaluer les risques de depassement
        risks = self._assess_budget_risks(pessimistic, tenant_id)
        
        return ForecastResult(
            optimistic=optimistic,
            baseline=baseline,
            pessimistic=pessimistic,
            risks=risks,
        )
\end{lstlisting}


\subsection{Pipeline Celery --- Tâches périodiques}

Les tâches d'agrégation et d'analyse sont orchestrées par \textbf{Celery Beat}~\cite{celery2024}, qui déclenche automatiquement les traitements selon un calendrier configuré~:

\begin{lstlisting}[language=Python3, caption={Configuration de Celery Beat (celery\_app.py)}]
app.conf.beat_schedule = {
    "aggregate-daily-costs": {
        "task": "modules.analytics.tasks.aggregate_daily_costs",
        "schedule": crontab(hour=2, minute=0),   # Chaque nuit a 2h
    },
    "rollup-monthly-costs": {
        "task": "modules.analytics.tasks.rollup_monthly_costs",
        "schedule": crontab(day_of_month=1, hour=3),  # 1er du mois
    },
    "run-optimization-scan": {
        "task": "modules.optimizer.tasks.run_weekly_scan",
        "schedule": crontab(day_of_week=1, hour=4),  # Chaque lundi
    },
}
\end{lstlisting}


\subsection{Captures d'écran}

\begin{figure}[H]
    \centering
    \fbox{\parbox{14cm}{\centering\vspace{4cm}\textit{Capture d'écran --- Page des anomalies détectées}\\[0.5cm]\textit{(Liste des anomalies avec sévérité, algorithme, date et tenant)}\\[0.5cm]\textit{(Détail d'une anomalie avec explication IA et recommandations)}\vspace{3cm}}}
    \caption{Interface de visualisation des anomalies détectées}
    \label{fig:anomalies_ui}
\end{figure}

\begin{figure}[H]
    \centering
    \fbox{\parbox{14cm}{\centering\vspace{4cm}\textit{Capture d'écran --- Page de prévision budgétaire}\\[0.5cm]\textit{(Graphique avec 3 scénarios~: optimiste (vert), moyen (bleu), pessimiste (rouge))}\\[0.5cm]\textit{(Tableau des risques de dépassement par tenant)}\vspace{3cm}}}
    \caption{Interface de prévision budgétaire à trois scénarios}
    \label{fig:forecasting_ui}
\end{figure}

\begin{figure}[H]
    \centering
    \fbox{\parbox{14cm}{\centering\vspace{4cm}\textit{Capture d'écran --- Page des coûts et analytics}\\[0.5cm]\textit{(Répartition des coûts par modèle, par provider, par jour)}\\[0.5cm]\textit{(Graphiques en barres et en aires empilées)}\vspace{3cm}}}
    \caption{Interface d'analyse des coûts}
    \label{fig:costs_ui}
\end{figure}


\subsection{Tests du Sprint~2}

\begin{table}[H]
\centering
\caption{Résultats des tests --- Sprint 2}
\label{tab:tests_s2}
\begin{tabular}{|l|c|c|c|}
\hline
\textbf{Module} & \textbf{Tests unitaires} & \textbf{Tests d'intégration} & \textbf{Résultat} \\
\hline
Detection & 18 & 6 & \textcolor{ForestGreen}{24/24 PASS} \\
\hline
Explainer & 10 & 4 & \textcolor{ForestGreen}{14/14 PASS} \\
\hline
Forecasting & 8 & 3 & \textcolor{ForestGreen}{11/11 PASS} \\
\hline
Analytics & 12 & 5 & \textcolor{ForestGreen}{17/17 PASS} \\
\hline
Optimizer & 7 & 3 & \textcolor{ForestGreen}{10/10 PASS} \\
\hline
\textbf{Total Sprint 2} & \textbf{55} & \textbf{21} & \textcolor{ForestGreen}{\textbf{76/76 PASS}} \\
\hline
\end{tabular}
\end{table}


% ────────────────────────────────────────────────────────────
\section{Sprint Rétrospective}
\label{sec:s2_retro}
% ────────────────────────────────────────────────────────────

\begin{table}[H]
\centering
\caption{Rétrospective du Sprint 2}
\label{tab:retro_s2}
\begin{tabularx}{\textwidth}{|l|X|}
\hline
\textbf{Ce qui a bien fonctionné} & \begin{minipage}[t]{\linewidth}\vspace{2pt}
\begin{itemize}[noitemsep,topsep=0pt]
\item La complémentarité des trois algorithmes de détection assure une couverture complète
\item Le bus d'événements permet un chaînage automatique~: détection $\rightarrow$ explication $\rightarrow$ notification
\item Le pipeline Celery gère efficacement les tâches d'agrégation nocturnes
\item L'intégration de pgvector permet la recherche sémantique d'anomalies similaires
\end{itemize}\vspace{2pt}
\end{minipage} \\
\hline
\textbf{Ce qui peut être amélioré} & \begin{minipage}[t]{\linewidth}\vspace{2pt}
\begin{itemize}[noitemsep,topsep=0pt]
\item Le seuil de détection de l'Isolation Forest nécessite un calibrage par tenant
\item Les appels LLM pour l'explication peuvent être coûteux~; le cache sémantique atténue ce problème
\end{itemize}\vspace{2pt}
\end{minipage} \\
\hline
\textbf{Actions pour le prochain sprint} & \begin{minipage}[t]{\linewidth}\vspace{2pt}
\begin{itemize}[noitemsep,topsep=0pt]
\item Développer le module de facturation contractuelle
\item Implémenter le traçage des sessions et la gestion des prompts
\item Mettre en place le déploiement Docker et le pipeline CI/CD
\end{itemize}\vspace{2pt}
\end{minipage} \\
\hline
\end{tabularx}
\end{table}


\section*{Conclusion}

Le Sprint~2 a doté la plateforme LLM Dashboard de capacités d'intelligence artificielle avancées. Le pipeline de détection d'anomalies, combinant STL, Isolation Forest et CUSUM, identifie de manière proactive les comportements de consommation anormaux. Le module Explainer produit des explications intelligibles et des recommandations actionnables grâce aux LLM. Le moteur de prévision budgétaire offre une visibilité prospective à trois scénarios, tandis que le pipeline Celery automatise l'ensemble des traitements d'agrégation. Les 76 tests validés confirment la fiabilité des implémentations. Le chapitre suivant détaille le Sprint~3, consacré aux fonctionnalités entreprise et au déploiement de la plateforme.
